% New Section 6. Editorial roles/evidence refer to DISCUSSION_WRITING_NOTES_ARS.md.
\section{讨论}
\label{sec:discussion}

\subsection{频谱细化与双耳线索的折衷}
% Role: interpretation of residual compensation. Evidence: E2 and Table primary-results.
MCA之后的残差学习在本研究中具有明确的频谱补偿价值。表\ref{tab:primary-results}及其配对结果显示，Hybrid降低了MCA的全域、对侧区域和对侧高频误差。这支持利用稀疏观测学习传统插值未消除的幅度残差。相对外部学习方法，收益则更具选择性：Hybrid在全域及对侧ERB上优于FSP-AE和RANF，但FSP-AE的对侧高频误差更低；相对MCA的水平面ILD收益也未获得配对区间支持。本文据此讨论不同误差目标间的取舍，不以领先项数量判定总体优劣。

% Role: endpoint interpretation and objective alignment; not a proven causal mechanism. Evidence: method/setup definitions, E1.
局部谱形与双耳能量关系反映不同的重建要求。一、二阶差分及多尺度凹口深度约束局部频谱变化，绝对幅度误差和ERB误差则评价其他尺度的偏离。Hybrid在三个局部端点上的最低均值，与FSP-AE的对侧高频优势可以同时成立。这些局部端点与训练辅助项定义相近，其改善说明模型在所约束的目标上表现较好，但不足以证明一般听觉收益。全域ERB和频带ILD在外部主比较中的最低均值，也不能保证严格水平面ILD的稳定收益；后者基于完整重建HRIR的双耳能量比。联合输入双耳信息并设置ILD损失，并不保证该误差在每名被试上下降。

\subsection{结构机制与模型容量}
% Role: design rationale versus isolated causal evidence. Evidence: S1/S2 and E4.
Hybrid将个体条件基场与方向条件化频谱卷积串联，使逐查询残差预测后仍可利用跨频率邻域及双耳信息进行修正。验证集组件比较中，追加并训练细化器后，三项主要频谱误差降低，而ILD差值区间含零（表\ref{tab:component-results}）。该观察支持完整细化方案的开发效果。由于比较同时引入CNN、额外训练和复合目标，尚不能把收益分别归因于卷积结构、某一损失项或集成。与训练目标一致的结构动机，应与已被隔离验证的机制区分。

% Role: initialization and efficiency boundary. Evidence: S1/S2, author-confirmed missing efficiency evaluation.
零输出初始化仅保证Hybrid的初始函数等于载入的FiLM父模型，不保证训练后的误差单调下降。冻结骨干减少了本阶段需要优化的参数，但完整推断仍包含骨干、CNN及三成员集成。因此，仅凭新增模块规模不能认定部署轻量或效率领先，其他内部变体的参数与效率结果也不能代替Hybrid的测量。

\subsection{泛化范围与研究局限}
% Role: scope of observation sensitivity. Evidence: E8 and Table q-results.
观测数量敏感性揭示了固定操作点的适用边界。在验证集共同743方向上，Q14的四项误差退化，Q50的全域与对侧ERB也未改善（表\ref{tab:q-results}）。这些变化可能与输入集合改变及固定模型的适配能力有关，但实验未分离其原因。结果仅适用于Q26训练并冻结后的响应，不能推出更多测量普遍有害，也不能替代针对新观测设置训练后的评价。

% Role: limits on generalization and confirmatory interpretation. Evidence: S3 and setup chronology.
证据范围还受单一数据库和开发历史限制。测试队列已有工程使用历史，Hybrid主叙事在结果已知后确定，不能视作新的未见独立确认。多端点比较未校正多重性，补充指标也不是相互独立的验证，均值接近和区间含零均不构成等效证据。本文结论限于所用被试队列、测量域及重建流程，跨数据库和其他测量条件下的适用性尚待检验。

% Role: individual failure and future work, not experiment authorization. Evidence: results retained case and author scope.
P0339的水平面ILD失败案例进一步表明，平均频谱收益不足以保证个体双耳线索质量。该被试保留在原统计中，失败原因尚未定位。后续研究优先针对水平面ILD与个体失败案例检验双耳约束及风险控制，再在独立数据上确认结果、隔离组件效应和测量完整推断成本。缺少听音实验时，客观谱误差及ITD合理性检查不能被解释为定位准确性、外化感或偏好的提升。这些验证均属未来工作，不计入本文贡献。
